\problemname{Mid Clocking}

We have a number of processor slots in a row. The processors are each dedicated to a single task
that they can complete before going into a busy-looping mode while they wait for
the next instruction.

A processor can run at 100\% speed if it has enough cooling, and that requires
there be no adjacent processor. If a processor has 1 adjacent processor, it only
runs at 60\% of maximum speed. If it has 2 adjacent processors it can only run at
30\% of maximum speed.

\begin{figure}[!h]
\centering
\includegraphics[height=0.15\textwidth]{sample2.pdf}
\caption{Illustration of Sample Input 2.}
\end{figure}

We know how many seconds each task will take running on a processor at full
speed, starting from $t=0$, and need to allocate tasks to processors to ensure
the last task finishes as soon as possible.

\section*{Input}

\begin{itemize}
  \item The number of tasks to run, $n$, and the number of processor slots, $s$ ($1 \le n \le s \le 10^6$).
  \item $n$ further integers giving the duration of each task $t_i$ in seconds if performed at 100\% speed ($1 \le t_i \le 10^9$).
\end{itemize}

\section*{Output}

Output the minimum time by which all of the tasks can be finished.

The output must be accurate to an absolute or relative error of $10^{-6}$.
